\section{Geometry of progress and stagnation}\label{art:geometry}

The mixed ridge $F(u,v)=u^2+|v|$ illustrates why a gradient jump need not
remove useful descent. Along $v=0$ its limiting gradients are $(2u,1)$ and
$(2u,-1)$. Their difference is vertical. At a point with $u\ne0$, the
horizontal vector $(-u,0)$ has the same strictly negative pairing $-2u^2$
with both gradients. The inward radial model also has a small error:
\[
 F((1-\tau)(u,v))-F(u,v)+\tau\nabla F(u,v)^T(u,v)
       =\tau^2u^2\qquad(v\ne0,\ 0\le\tau\le1).
\]
Thus the nonsmooth part is exactly linear under inward scaling. This is the
mechanism formalized by the radial theorem, not an assertion about the eventual
eigenvectors or eigenvalues of a computed BFGS matrix.

\begin{figure}[htbp]
\centering
% Vector construction from explicit gradients; no trajectory inference.
\begin{tikzpicture}[x=1.7cm,y=1.5cm,>=Stealth,font=\small]
\begin{scope}[shift={(-1.9,0)}]
\draw[->,gray] (-.8,0)--(1.55,0);\draw[->,gray] (0,-1.35)--(0,1.35);
\draw[line width=1.1pt] (1,-1)--(1,1);
\draw[->] (0,0)--(1,1);\draw[->] (0,0)--(1,-1);
\fill (1,1) circle (1.4pt) node[right] {$(1,1)$};
\fill (1,-1) circle (1.4pt) node[right] {$(1,-1)$};
\draw[->,thick,dashed] (0,0)--(-.65,0) node[below] {descent};
\node at (.4,-1.65) {(a) $u^2+|v|$, $u=1/2$};
\node[above left] at (0,0) {$0$};
\end{scope}
\begin{scope}[shift={(1.8,0)}]
\draw[->,gray] (-1.35,0)--(1.4,0);\draw[->,gray] (0,-1.35)--(0,.55);
\fill[gray!15] (1,0)--(-1,-.125)--(.3184814276,-1.0252321197)--cycle;
\draw[thick] (1,0)--(-1,-.125)--(.3184814276,-1.0252321197)--cycle;
\draw[->] (0,0)--(.3184814276,-1.0252321197);
\fill (1,0) circle (1.4pt) node[above right] {$g_0$};
\fill (-1,-.125) circle (1.4pt) node[below left] {$g_1$};
\fill (.3184814276,-1.0252321197) circle (1.4pt) node[right] {$g_2$};
\fill[white] (0,0) circle (1.7pt);\draw (0,0) circle (1.7pt);
\node[above] at (0,.04) {$0$};
\node at (0,-1.65) {(b) Certified nonstationary seed};
\end{scope}
\end{tikzpicture}

\caption{Gradient space. Left: the mixed ridge at $u=1/2$ has a vertical
jump and a common horizontal descent direction. Right: the particular
three-gradient seed used in Section~\ref{art:failure}. Its convex hull
excludes zero, but its affine hull is the plane. The vertices on the right
are rendered from the certified root's rational center; the mathematical
definition uses the isolated root. Arrows indicate gradients, not BFGS steps.}
\label{art:gradient-geometry}
\end{figure}

The bad corner in Figure~\ref{art:gradient-geometry} is specific. Its slopes
are $g_0=(1,0)$, $g_1=(-1,-1/8)$ and $g_2=(c,d)$, where $(c,d)$ is
defined exactly in Appendix~\ref{art:seed_certificate}. They are affinely
independent, so no nonzero direction pairs equally with all three. Yet
$(-1,16)$ has negative pairing with each, and zero is outside their convex
hull. The corner is therefore nonstationary. Section~\ref{art:failure}
shows that the actual algorithm nevertheless contracts through a three-step
cycle toward this corner, with all line-search comparisons strict.

This distinction between affine and convex balance helps interpret the
examples. For a nonempty finite family of slopes $a_i$, equality $a_i^Tv=\kappa<0$ for every $i$ asks
for a common rate of descent, whereas $a_i^Tv<0$ asks only for descent.
The former fails when zero belongs to their affine hull; the latter fails
when zero belongs to their convex hull. Indeed, a linear functional constant
on the slopes is constant on their affine hull. Conversely, if zero is
outside that affine hull, its nearest point gives a separating functional
constant on the hull; changing its sign gives the negative common value.
The convex version follows from the same closest-point argument for the
convex hull. These are elementary separation facts, not a classification
of BFGS behavior.

At the minimizer of the mixed ridge there is no remaining objective gap to
eliminate. At the certified wrong corner, contraction occurs while a strict
descent direction remains; its coercive realization in
Corollary~\ref{art:failure_polyhedral} also has a positive limiting objective gap.
A finite number of smooth
formulas can repeat the latter configuration at infinitely many shrinking
scales. The local positive result instead controls that repetition by a finite
polyhedral partition near the optimizer. The proof makes this control
quantitative through the error in an inward first-order model.
